\documentclass[11pt]{article}
\usepackage[a4paper,margin=27mm,headheight=14pt]{geometry}
\usepackage[T1]{fontenc}
\usepackage{lmodern,microtype,amsmath,amssymb,amsthm,mathtools,booktabs,array}
\usepackage[dvipsnames]{xcolor}
\usepackage[colorlinks=true,linkcolor=MidnightBlue,citecolor=MidnightBlue,urlcolor=MidnightBlue]{hyperref}
\hypersetup{pdftitle={Exact residue transport for sparse p-adic roots},pdfsubject={Sharp geometric root counts and finite optimization}}
\usepackage{fancyhdr}
\pagestyle{fancy}\fancyhf{}\lhead{\small\textsc{Sparse $p$-adic root counts}}\rhead{\small Research note}\cfoot{\thepage}
\newtheorem{theorem}{Theorem}[section]
\newtheorem{lemma}[theorem]{Lemma}
\newtheorem{corollary}[theorem]{Corollary}
\newtheorem{proposition}[theorem]{Proposition}
\theoremstyle{remark}\newtheorem{remark}[theorem]{Remark}
\newcommand{\Q}{\mathbb Q}\newcommand{\Z}{\mathbb Z}\newcommand{\F}{\mathbb F}\newcommand{\C}{\mathbb C}
\newcommand{\vp}{v_p}\newcommand{\ord}{\operatorname{ord}}\newcommand{\spn}{\operatorname{span}}\newcommand{\res}{\operatorname{res}}
\newcommand{\D}{\mathcal D}\newcommand{\R}{\mathcal R}\newcommand{\Ocal}{\mathcal O}
\setlength{\emergencystretch}{2em}
\title{Exact residue transport for sparse $p$-adic roots\\[3pt]\large Arbitrary supports, exponential polynomials, and the range from $p+2$ to $2p+1$ terms}
\author{Research note prepared for the ProveIt project}
\date{27 November 2026}
\begin{document}
\maketitle
\begin{abstract}
For a prime $p$ and $t\ge2$, write $S_p(j)$ for the base-$p$ digit sum, set
$D=\max_{0\le j<t}S_p(j)$ and $A=\lceil(D-1)/(p-1)\rceil$.
We prove that the complete multiset of depths of the roots of a $t$-term polynomial in
$v_p(x-1)>1/(p-1)$ can be transported, in both directions, between any two supports with
identical exponent multisets modulo $p^A$. In fact, this extends to exponential
polynomials with frequencies in $\Z_p$ and polynomial amplitudes, counting frequencies
with their differential multiplicities. Thus every fixed-support extremum at every
closed depth, as well as the universal open-disk extremum, reduces exactly to a finite
integer matrix problem, with representative polynomial degree at most $tp^A-1<pt^2$.
The method determines the exact extremum for all primes in the range
\[
 p+2\le t\le 2p+1:\qquad N_p^\circ(t)=p^2+\left\lfloor\frac{t-p}{2}\right\rfloor-1.
\]
The proof gives a uniform two-residue construction attaining the bound with distinct
roots and integer coefficients. These results extend the residue-algebra methods of the
ProveIt GRH article beyond its previous restriction on support size. All results here
are unconditional; none relies on RH, Dirichlet GRH, or Artin holomorphy.
\end{abstract}

\section{Problem, scope, and main statements}

The ProveIt article \cite{article} asks, in its question \texttt{smd:q:roots}, which additional
residue information controls the sharp geometric root bound beyond $p$ terms. Its
results \texttt{smd:thm:rankrule}, \texttt{smd:thm:fixedsupport} and
\texttt{smd:thm:opendisk} determine the answer for supports of size at most $p$; its
\texttt{smd:thm:jump} reaches $p+1$ and $p+2$ terms and in particular proves
$N_p^\circ(p+2)=p^2$. The obstruction to simply repeating the mod-$p$ proof is that
$T^p$ may have a nonzero nilpotent part in the residue algebra. We replace the
single-digit argument by an integral moment comparison with a controlled congruence
precision, then obtain a finite exact optimization by a determinant argument. The
closed-form extension exploits the residual nilpotent part rather than discarding it.

Let $p$ be prime, normalize $\vp(p)=1$, and work over $\overline{\Q}_p$ (or its completion
$\C_p$). For $a\ne0$ and real $r>r_0:=1/(p-1)$, set
\[
 \D(a,r)=\{x:\vp(x/a-1)\ge r\},\qquad
 \D^\circ(a)=\{x:\vp(x/a-1)>r_0\}.
\]
For a finite set $E$ of distinct nonnegative integers, $N_p(E,r)$ is the largest
number of distinct roots in $\D(1,r)$ of a nonzero polynomial supported on $E$;
coefficients may vanish and lie in any finite extension of $\Q_p$.
Define $N_p^\circ(E)$ similarly for $\D^\circ(1)$, and $N_p(t,r),N_p^\circ(t)$ by
allowing at most $t$ terms. Scaling by $a$ does not change these maxima. We prove
upper bounds counting multiplicities; below we show that the maxima can be attained
with all roots in the disk simple. The depth of the center itself is $+\infty$.

For an integer $t\ge2$, define
\begin{equation}\label{eq:precision}
 D_{p,t}=\max_{0\le j<t} S_p(j),\qquad
 A_{p,t}=\left\lceil\frac{D_{p,t}-1}{p-1}\right\rceil,\qquad M_{p,t}=p^{A_{p,t}}.
\end{equation}
Here $D_{p,t}\ge1$ and $A_{p,t}\ge0$; in particular modulus $1$ is allowed.
For $t=1$ put $M_{p,1}=1$ and all maxima equal zero.

\begin{theorem}[Exact residue transport]\label{thm:transport}
Suppose $E,E'$ have the same size $t$ and give the same multiset modulo $M_{p,t}$.
There is a bijective linear coefficient transport between their polynomial spaces
(after fixing the center) that preserves the multiset of depths of all roots in
$\D^\circ(1)$, including multiplicities. In particular, using simple-root attainment (proved below), for every $r>r_0$,
\[
 N_p(E,r)=N_p(E',r),\qquad N_p^\circ(E)=N_p^\circ(E').
\]
More generally, consider a nonzero exponential polynomial
\begin{equation}\label{eq:exp}
 f(z)=\sum_{i=1}^{s}\sum_{j=0}^{m_i-1} c_{ij}z^j\exp(\lambda_i z),\qquad
 \lambda_i\in\Z_p\text{ distinct},\quad m_i\ge1,\quad \sum_i m_i=t,
\end{equation}
on $\vp(z)>r_0$. If $E'$ consists of $t$ distinct nonnegative integers with the same
residue multiset modulo $M_{p,t}$ as the $\lambda_i$ repeated $m_i$ times, then the
root-depth multiset of $f$ in the stated disk agrees with that of some
nonzero polynomial $P$ supported on $E'$ in $\D^\circ(1)$, and conversely every such
polynomial's multiset is attained in the exponential-polynomial space.
\end{theorem}

\begin{theorem}[Finite sharp determination]\label{thm:algorithm}
For $t=|E|\ge2$, let $\nu_c$ be the multiplicity of residue $c$ in $E$ modulo $M=M_{p,t}$ and define
\begin{equation}\label{eq:representative}
 E^*=\bigcup_{c=0}^{M-1}\{c+Mj:0\le j<\nu_c\},\qquad d=\max E^*.
\end{equation}
Use the $d+1$ rows $B_j=(\binom{e}{j})_{e\in E^*}$, $0\le j\le d$, with $\binom ej=0$ if $j>e$.
For each $t$-element index set $I\subseteq\{0,\ldots,d\}$ with nonzero determinant
$\Delta_I=\det(B_j)_{j\in I}$, define
\[
 W_r(I)=\vp(\Delta_I)+r\sum_{j\in I}j.
\]
Then
\begin{equation}\label{eq:matrix}
 N_p(E,r)=\max\{\max I: I\text{ minimizes }W_r\}.
\end{equation}
For $N_p^\circ(E)$ use the same formula but minimize the pair
\begin{equation}\label{eq:lex}
 \left((p-1)\vp(\Delta_I)+\sum_{j\in I}j,\ \sum_{j\in I}j\right)
\end{equation}
lexicographically. Maximizing over the finitely many size-$t$ residue multisets
modulo $M$ yields $N_p(t,r)$ and $N_p^\circ(t)$ exactly. In particular,
\begin{equation}\label{eq:degreebound}
 N_p(E,r)\le N_p^\circ(E)\le\max E^*-\min E^*\le tM-1<pt^2.
\end{equation}
The maximum at every depth is attained by an integer polynomial with every exponent
of the original support present and all nonzero roots simple, which becomes irreducible over
$\Q$ after division by its smallest monomial.
\end{theorem}

\begin{theorem}[Sharp new range]\label{thm:sharp}
For every prime $p$ and $p+2\le t\le2p+1$,
\begin{equation}\label{eq:sharp}
 \boxed{\displaystyle N_p^\circ(t)=p^2+\left\lfloor\frac{t-p}{2}\right\rfloor-1.}
\end{equation}
The bound counts roots with multiplicity and is attained by an integer polynomial with
simple roots. Writing $h=\lfloor(t-p)/2\rfloor$ and $K=p^2+h-1$, a witness uses the support
\begin{equation}\label{eq:witnesssupport}
 E_h=\{0,p,2p,\ldots,(p+h-1)p\}\ \cup\ \{1,1+p,\ldots,1+(h-1)p\}
\end{equation}
and has $K$ distinct roots already in any closed disk of rational depth
$r=r_0+\delta$ with $0<\delta K<1/(p-1)$, with suitable integral coefficients.
\end{theorem}

Theorem~\ref{thm:transport} is stronger than just invariance of the extremum: the
same coefficient transport works simultaneously for every depth. Theorem~\ref{thm:algorithm}
answers the question of finite determination for arbitrary $t$ and for deeper closed
disks; the number of residue profiles may be large, so it is not a polynomial-time
claim. Theorem~\ref{thm:sharp} is an explicit infinite family of new exact values,
including $t=p+3$ through $2p+1$ (where nonempty), beyond the known endpoint $p+2$.

\section{Integral moments and congruence precision}

We give the argument in the more general exponential-polynomial formulation. Let
$g(T)=\prod_{i=1}^s(T-\lambda_i)^{m_i}\in\Z_p[T]$ and
$\Lambda=\Z_p[T]/g(T)$. This is a free $\Z_p$-module with basis $1,T,\ldots,T^{t-1}$.
For the polynomial case put $\lambda_i=e_i\in E$, $m_i=1$, and apply the proof to
$f(z)=P(\exp z)$.

The derivatives at the origin satisfy
\begin{equation}\label{eq:moments}
 M_k=f^{(k)}(0)=\sum_{i,j}c_{ij}(T^k)^{(j)}(\lambda_i)=L(T^k),\qquad
 L(u)=\sum_{i,j}c_{ij}u^{(j)}(\lambda_i).
\end{equation}
Every derivative of order $j<m_i$ annihilates multiples of $g$ at $\lambda_i$, so $L$
is a well-defined functional on the scalar extension of $\Lambda$. Over a field of
characteristic zero the Chinese remainder decomposition into the powers of the distinct
linear factors gives Hermite interpolation. In each factor the derivatives of orders
$0,\ldots,m_i-1$ give independent coordinates (the factors $j!$ are nonzero).
Consequently the $t$ values $M_0,\ldots,M_{t-1}$ are arbitrary and determine the
coefficients $c_{ij}$ uniquely. In particular a nonzero $f$ has some nonzero initial
moment. Division by a nonzero initial moment of least valuation normalizes $L$ so that
\begin{equation}\label{eq:normal}
 L(1),L(T),\ldots,L(T^{t-1})\in\Ocal,
 \qquad \vp(M_{i_0})=0\quad\text{for some }i_0<t,
\end{equation}
where $\Ocal$ is the valuation ring of a field containing the coefficients. Since $g$
is monic integral, the coordinates of every $T^k$ in this power basis are integral.
Thus every normalized $M_k$ has nonnegative valuation. Equivalently the moments satisfy
the integral monic recurrence defined by $g$.

For rational $r>r_0$, choose $b\in\overline{\Q}_p$ with $\vp(b)=r$. The power series
$f(bY)=\sum_{k\ge0}M_k b^kY^k/k!$ converges on the closed unit disk. Its coefficients
have valuations
\begin{equation}\label{eq:weights}
 \vp(M_k)+w_r(k),\qquad
 w_r(k)=rk-\vp(k!)=(r-r_0)k+\frac{S_p(k)}{p-1}.
\end{equation}
Legendre's formula was used in the last equality. Since $r-r_0>0$, the valuations
tend to infinity. By the Newton polygon theorem, equivalently Weierstrass preparation
for a convergent one-variable restricted series, its number of zeros in the closed unit
disk over $\C_p$, with multiplicities, is the largest index realizing the minimum of
these valuations. Under $x=\exp z$ we have $\vp(x-1)=\vp(z)$, and exp and log are inverse
analytic bijections on these disks, with nonvanishing derivatives. Thus this also counts
polynomial roots on $\D(1,r)$ in the polynomial case.

\begin{lemma}[Excluding high-valuation moments]\label{lem:exclude}
Under \eqref{eq:normal}, let $k\ge t$. If $\vp(M_k)\ge A=A_{p,t}$, then the coefficient
of index $k$ cannot realize the minimum in \eqref{eq:weights}. If $k\ge t$ does realize
that minimum, then
\[
 S_p(k)\le D_{p,t}-1,\qquad \vp(M_k)<A.
\]
\end{lemma}
\begin{proof}
Put $\delta=r-r_0>0$. Since $k>i_0$ and $S_p(k)\ge1$,
\begin{align*}
 w_r(i_0)-w_r(k)
 &=\delta(i_0-k)+\frac{S_p(i_0)-S_p(k)}{p-1}\\
 &<\frac{D_{p,t}-S_p(k)}{p-1}
 \le\frac{D_{p,t}-1}{p-1}\le A.
\end{align*}
So $\vp(M_k)\ge A$ gives $\vp(M_k)+w_r(k)>w_r(i_0)$, strictly excluding $k$.
If $S_p(k)\ge D_{p,t}$, then the first strict inequality gives
$w_r(i_0)-w_r(k)<0$, and $\vp(M_k)\ge0$ again excludes it. This also treats
$A=0$: all indices $k\ge t$ are then excluded.
\end{proof}

\begin{proof}[Proof of Theorem~\ref{thm:transport}]
Let $g'$ be another monic split polynomial of the same degree whose roots with
multiplicity are congruent as a multiset modulo $p^A$, and use the corresponding
moment functional $L'$ with precisely the same first $t$ moments as $L$.
Hermite interpolation shows that this matching is a linear bijection before any
normalization. Since $g\equiv g'\pmod{p^A}$, their monic remainders of every $T^k$
have coordinatewise congruent coefficients in the respective power bases. After
common normalization as in \eqref{eq:normal}, it follows that
\begin{equation}\label{eq:congmoment}
 M_k\equiv M'_k\pmod{p^A\Ocal}\quad(k\ge0).
\end{equation}
For $k<t$ the moments coincide exactly. For $k\ge t$ there are two possibilities.
If one of the moments has valuation less than $A$, the other has exactly the same
valuation by \eqref{eq:congmoment}. Otherwise both have valuation at least $A$, and
Lemma~\ref{lem:exclude} excludes both from the minimum. Hence for every rational
$r>r_0$ the sets of minimizing indices in the two series coincide. The root count
at every such closed depth, with multiplicity, is identical.

In particular, choose $g'$ from any distinct integer representatives of the
residue multiset. Its exponential polynomial is $P(\exp z)$ for an ordinary
polynomial $P$. Its total number of roots in the open disk is finite. Therefore the
original exponential polynomial also has finitely many roots there: any finite
collection of roots would lie in a single closed disk of rational depth $>r_0$.
The common counts at all rational depths determine the multiset of depths, including
$+\infty$, by taking differences and limits. This proves the theorem for the entire
open disk and for all real closed depths, not just rational ones. No hypothesis on
the simplicity of the original roots is needed.
\end{proof}

\begin{remark}[Why ramification is harmless]
The congruence is modulo $p^A\Ocal$, not the maximal ideal of a potentially highly
ramified coefficient field. A valuation less than $A$ may be fractional, but the
ultrametric inequality still gives equality of valuations under
\eqref{eq:congmoment}. This distinction is necessary for geometric root counts
rather than counts restricted to $\Q_p$.
\end{remark}

\section{Sharp finite optimization at every depth}

We prove Theorem~\ref{thm:algorithm} and describe a practical exact basis-exchange
algorithm. The support $E^*$ in \eqref{eq:representative} is valid by
Theorem~\ref{thm:transport}. In particular the degree of its polynomial after
removing the smallest monomial is at most $tM-1$. If $t\ge2$ and
$H=\lceil\log_p t\rceil$, then $D_{p,t}\le(p-1)H$, so $A\le H$ and
$M\le p^H<pt$. This proves \eqref{eq:degreebound}. (When $t=1$, the stated bound on
$N$ is interpreted directly.) The same argument applies to any exponential-polynomial
space in Theorem~\ref{thm:transport}, without assuming the existence of an extremum
in advance.

\begin{lemma}[A minimum-determinant basis detects all possible last minima]\label{lem:lattice}
Let $u_0,\ldots,u_d$ be row vectors spanning $K^t$ over a discretely valued field $K$
with valuation ring $\Ocal_K$, and let $\mathcal L=\sum_j\Ocal_K u_j$. Among all bases
chosen from these rows, take one of minimum determinant valuation. Its rows form an
$\Ocal_K$-basis of $\mathcal L$. A row belongs to at least one minimizing basis if and
only if it has a unit coordinate in this basis. The largest index that can realize
the minimum valuation of the sequence $u_j c$ for a nonzero column $c\in K^t$ is
therefore the largest index belonging to some minimizing basis.
\end{lemma}
\begin{proof}
Write the coordinates of a row $u_j$ in a minimizing basis as $a_i$.
Cramer's rule expresses $a_i$ as a determinant ratio. A coordinate of negative
valuation would give a smaller determinant by replacing row $i$, which is impossible.
Thus all coordinates are integral and the selected basis generates $\mathcal L$.
A unit coordinate yields a basis exchange of equal determinant valuation.
Conversely, if all coordinates lie in the maximal ideal, so does the row in
$\mathcal L$; it cannot occur in any basis of $\mathcal L$.
For a nonzero column $c$, the smallest valuation on all rows equals its smallest
valuation on basis rows, and a row with all coordinates nonunits gives a strictly
larger valuation. If a row belongs to a minimum basis, choose a dual column that
vanishes on every other row in that basis and takes value one on it. Every row then
has value integral, so the selected row realizes the minimum. Taking the largest
index that can do this proves the last assertion.
\end{proof}

\begin{proof}[Proof of the finite formulas in Theorem~\ref{thm:algorithm}]
For rational $r>r_0$ choose $b$ of valuation $r$ and write
\[
 P(1+bY)=\sum_{j=0}^d b^j(B_j c)Y^j.
\]
The $t$ columns are independent, since the polynomials $(1+bY)^e$ for distinct $e$
are independent. Apply Lemma~\ref{lem:lattice} to $u_j=b^j B_j$. A basis determinant
has valuation $W_r(I)$. The polynomial Newton polygon theorem now gives
\eqref{eq:matrix}. The same proof works if the coefficient field is enlarged, since
all the coordinate valuations and their comparisons remain unchanged. Although the
basis argument used a field containing $b$, the attaining coefficient column can be
chosen rational: for a minimizing basis containing the largest possible index $k$,
take the rational column satisfying $B_kc=1$ and $B_jc=0$ on the other basis rows.
Then the valuations of every $b^j B_jc$ are at least $rk$ by Cramer's rule, and the
$k$th realizes equality. The largest possible minimizing index is exactly $k$.

For a general real $r$, only finitely many affine functions $W_r(I)$ occur. The
crossings of unequal such functions are rational, since the determinant valuations
and index sums are integers. At an irrational $r$ choose rational depths on both
sides sufficiently close that the minimizing bases do not change. Monotonicity in
$r$ proves the formula there as well. At rational crossing points the use of all
minimizers in \eqref{eq:matrix} accounts exactly for roots on the boundary.

Write $r=r_0+\delta$. If $0<\delta<1/((p-1)t(d+1))$, the order of the determinant
scores is the lexicographic order of \eqref{eq:lex}: any distinct first entries differ
by at least one, while the possible change in $(p-1)\delta\sum j$ is less than one.
Therefore the maximum is constant for every sufficiently small positive $\delta$.
Any finite collection of roots in the open disk is contained in such a closed disk
of sufficiently small positive $\delta$, which proves the open formula.

For any polynomial of at most $t$ terms, pad its support with distinct unused
exponents up to size $t$. All possible supports of size $t$ have one of the finitely
many residue multisets and every multiset has the representative \eqref{eq:representative}.
Thus maximizing the exact formula over all these multisets gives the universal
maxima, with no relaxation. The case $t=1$ is immediate.
\end{proof}

\begin{proof}[Attainment and rational irreducible witnesses]
The rational coefficient construction works on the original support $E$ too, by
applying the same determinant argument to the Taylor matrix of $E$ itself (finite,
though possibly large); Theorem~\ref{thm:transport} gives its value from $E^*$.
For irrational depth, use a nearby rational depth on the deeper side with equal
maximal value. For the open disk use a sufficiently shallow rational depth as above.
Scale the resulting polynomial to integer coefficients and divide by its smallest
allowed monomial (even if the corresponding coefficient is currently zero).
Let $\mu$ be the minimum valuation of the Taylor coefficients after scaling the
chosen disk to the unit disk. Changing any integer coefficient by a multiple of
$p^B$, with integer $B>\mu$, changes every rescaled Taylor coefficient by a term of
valuation at least $B$, since binomial coefficients are integral and $r>0$.
Hence the minimum and its last index are unchanged.

Choose a prime $\ell\ne p$. By the Chinese remainder theorem we can independently
choose the changes in the allowed integer coefficients so that the leading coefficient
is nonzero mod $\ell$, all other allowed coefficients are zero mod $\ell$, and the
constant coefficient has valuation exactly one at $\ell$. Absent coefficients stay
zero. Choose all allowed coefficients nonzero (each congruence class is infinite).
The resulting polynomial is Eisenstein at $\ell$, hence irreducible over $\Q$ and
separable in characteristic zero. Multiplying back the smallest monomial introduces
no roots in the multiplicative disk. This gives the promised distinct roots with
exactly the extremal count.
\end{proof}

\paragraph{Basis exchanges rather than enumerating minors.}
The first $t$ Taylor rows are independent: evaluating polynomials in $e$ of degrees
$0,\ldots,t-1$ with nonzero leading coefficients at distinct $e$ gives an invertible
matrix. Start from this basis. When adding a row $B_j$, find its coordinates $a_i$ in
the current \emph{unscaled} basis with index list $I=(i_1,\ldots,i_t)$. In the scaled
basis the coordinate valuations are $\vp(a_i)+r(j-i_i)$. If all are nonnegative, the
row is already in the lattice. Otherwise choose a coordinate of minimum valuation
and replace that basis row. Solving for the old row in terms of the new rows shows that
all coefficients are integral, since the pivot has minimum valuation. Hence the new
basis spans precisely the enlarged lattice. After all rows are processed, the largest
row with minimum coordinate valuation zero gives the answer. This avoids exhaustive
enumeration of $\binom{d+1}{t}$ minors for a fixed profile. The accompanying Python
program implements this algorithm with rational arithmetic and the open-disk choice
$\delta=1/[2(p-1)t(d+1)]$.

\section{Proof of the sharp formula beyond $p+2$ terms}

Let $p+2\le t\le2p+1$, and keep the notation $h=\lfloor(t-p)/2\rfloor$.
For $p=2$ this means $t=4,5$ and $h=1$. For $p\ge3$ we have $1\le h<p$.
All the initial digit sums $S_p(i)$ for $0\le i<t$ are at most $p$, including the
case $p=2$. Thus $A_{p,t}\le1$; in fact $A_{p,t}=1$ in this range. For a normalized
moment sequence any minimizing index $k\ge t$ satisfies
\begin{equation}\label{eq:eligible}
 S_p(k)\le p-1,\qquad \vp(M_k)<1
\end{equation}
by Lemma~\ref{lem:exclude}. A support with fewer than $t$ terms may be padded to size
$t$ throughout the upper-bound proof.

\begin{lemma}[Reduction test]\label{lem:reduction}
Let $\R=\F_p[T]/\prod_{e\in E}(T-\bar e)$. If $k\ge t$ satisfies
\eqref{eq:eligible}, and in $\R$ the class of $T^k$ is a linear combination of classes
$T^j$ with $j<k$ and $S_p(j)\le S_p(k)$, then $k$ cannot minimize the coefficient
valuations at any $r>r_0$.
\end{lemma}
\begin{proof}
Lift the relation into the integral power basis. The moments satisfy the same
relation modulo $p\Ocal$ by integrality of $L$. As $\vp(M_k)<1$, at least one moment
on the right has valuation at most $\vp(M_k)$. Every index on the right has
strictly smaller weight by \eqref{eq:weights}. Hence some coefficient has strictly
smaller valuation, as required.
\end{proof}

\subsection{Upper bound: at most one thick residue factor}

For clarity we only have to exclude indices $k>p^2+h-1$; all such indices are $\ge t$,
since $t-1\le2p\le p^2$. Assume that the last minimizing index $k$ at a rational depth $r>r_0$ satisfies $k>p^2+h-1$, and let
\[
 a=k\bmod p,\qquad b=\sum_{u\ge1}a_u,\qquad k=\sum_{u\ge0}a_up^u.
\]
By \eqref{eq:eligible}, $a+b\le p-1$ and $b\ge1$.
The compressed index $j=a+pb$ has digit sum $a+b$ and satisfies $j<p^2<k$.

If all support exponents belong to one residue class modulo $p$, the exact transport
of Theorem~\ref{thm:transport} replaces the support by $c,c+p,\ldots,c+p(t-1)$.
After removing a monomial, such a polynomial has the form $Q(X^p)$ with
$\deg Q\le t-1$. The map $x\mapsto x^p$ is injective on $\D^\circ(1)$ because it
corresponds to multiplication by $p$ under the logarithm, and its derivative does not
vanish. Hence there are at most $t-1$ roots, with multiplicity. We may now suppose
that at least two residues occur.

If every residue multiplicity is at most $p$, the Chinese remainder factors of $\R$
are $\F_p[T]/(T-c)^{\nu_c}$. On each such factor $T^p=c$ and $T^{p^u}=c$ for every
$u\ge1$. Hence $T^k=T^{a+pb}$, contradicting Lemma~\ref{lem:reduction}.

Otherwise exactly one residue has multiplicity exceeding $p$, because $t\le2p+1$.
Multiplying the original polynomial by a nonnegative monomial, which preserves its
roots and multiplicities in the disk, shifts this residue to $0$; apply the normalized
moment argument to the shifted support. Since the root count is unchanged, its last
minimizing index is still $k$, by the Newton polygon theorem. Let the heavy multiplicity be $p+s$ and the
sum of all other multiplicities be $m$. Then
\begin{equation}\label{eq:profile}
 1\le s\le p,\qquad m\ge1,\qquad p+s+m=t,
 \qquad \R\cong\F_p[T]/T^{p+s}\times\R_{\rm other}.
\end{equation}
Every factor in $\R_{\rm other}$ has multiplicity at most $p$. On these factors,
$T^k=T^{a+pb}$ as before. On the zero-residue factor $T^k=0$, since $k>p^2\ge p+s$.
If $a+pb\ge p+s$ then the equality holds on the zero factor too, and the reduction
test excludes $k$. Since $p+s\le2p$, this covers $b\ge2$ and also $b=1,a\ge s$.

The only remaining possibility is $b=1,a<s$. Therefore $k=p^u+a$ for some $u\ge2$,
with $a\le p-2$. For every such $u$, $T^{p^u}=T^{p^2}$ in $\R$: both vanish on the
zero factor and both are the constant residue on each other factor. If $u>2$,
compress to $p^2+a$ and use Lemma~\ref{lem:reduction}. Thus the only possible new
index above $p^2$ is $k=p^2+a$, with $a<s$.

On $\R_{\rm other}$ the product of its defining primary polynomials is monic of
degree $m$ and annihilates $T$. On the zero factor $T^{p^2}=0$. Consequently for
$a\ge m$, the class $T^{p^2+a}$ is a linear combination of the classes
$T^{p^2+j}$, $0\le j<m$. Each has smaller index and digit sum $1+j<1+a$.
Lemma~\ref{lem:reduction} excludes these $a$ as well. A surviving index must obey
\[
 a\le\min(s,m)-1\le\left\lfloor\frac{s+m}{2}\right\rfloor-1=h-1.
\]
This contradicts the assumed $k>p^2+h-1$. The argument is valid for every rational
$r>r_0$ and hence throughout the open disk, since each finite set of roots is contained
in such a closed disk. It proves the upper bound even with multiplicities.

\subsection{Matching lower bound: an explicit functional on two residue factors}

Use the support $E_h$ of \eqref{eq:witnesssupport}, of size $t_0=p+2h\le t$.
Its residue algebra is
\[
 \R_h=\F_p[T]/T^{p+h}\times\F_p[T]/(T-1)^h.
\]
In the second factor put $Z=T-1$. On a pair $(u(T),v(1+Z))$ in this algebra define
\begin{equation}\label{eq:lambda}
 \lambda(u,v)=[Z^{h-1}]v(1+Z)-[T^{h-1}]u(T)
 -h[T^h]u(T)-[T^{p+h-1}]u(T),
\end{equation}
where each coefficient is taken in the unique truncated representative in its local
variable. For an index $j\ge0$, its value on $T^j$ is thus
\begin{equation}\label{eq:lambdaexplicit}
 \lambda(T^j)=\binom{j}{h-1}-\mathbf1_{j=h-1}
 -h\mathbf1_{j=h}-\mathbf1_{j=p+h-1}\quad\text{in }\F_p.
\end{equation}
This includes $h=1$ without modification. Define $K=p^2+h-1$.

\begin{lemma}[The decisive moment]\label{lem:decisive}
One has $\lambda(1)=0$, $\lambda(T^K)=1$, and
\[
 \lambda(T^j)=0\quad\text{if}\quad
 S_p(j)<h\quad\text{or}\quad(S_p(j)=h\text{ and }j<K).
\]
\end{lemma}
\begin{proof}
Since $h-1<p$, the binomial coefficient in \eqref{eq:lambdaexplicit} is congruent
to $\binom{a}{h-1}$, where $a=j\bmod p$. This follows by reducing
$(1+Z)^j=(1+Z)^a\prod_{u\ge1}(1+Z^{p^u})^{a_u}$ in characteristic $p$.
If $a<h-1$, the binomial coefficient and all three correction terms vanish.
If $a\ge h-1$ and $S_p(j)\le h$, the only possibilities are
$j=h-1$, $j=h$, and $j=p^u+h-1$ for $u\ge1$.
Among these, the indices with smaller digit sum or equal digit sum and $j<K$ are
exactly $h-1,h,p+h-1$, for which the three correction terms cancel the binomial
coefficients $1,h,1$, respectively. At $K$ none of the corrections occurs and the
binomial coefficient is one mod $p$. The assertion for $1=T^0$ follows from the
same computation (also directly from \eqref{eq:lambdaexplicit}).
\end{proof}

Lift $\lambda$ to an integral functional $L$ on
$\Z[T]/\prod_{e\in E_h}(T-e)$ by lifting its values on the power basis to integers,
choosing $L(1)=0$ exactly. The moments $M_j=L(T^j)$ are integral and reduce to
$\lambda(T^j)$ modulo $p$. By Lagrange interpolation there are rational coefficients
$c_e$ such that $M_j=\sum_{e\in E_h}c_e e^j$ for all $j$. Explicitly, with
\[
 \ell_e(T)=\prod_{f\in E_h,\ f\ne e}\frac{T-f}{e-f},
 \qquad c_e=L(\ell_e),
\]
where $L$ is extended rationally, this follows from interpolation modulo the defining
polynomial. It produces a nonzero polynomial $P(X)=\sum c_e X^e$.

Take $r=r_0+\delta$ rational with $0<\delta K<1/(p-1)$. Since $S_p(K)=h$,
\[
 w_r(K)=\frac h{p-1}+\delta K<\frac{h+1}{p-1}\le\frac p{p-1}.
\]
If $w_r(j)<w_r(K)$, then $S_p(j)\le h$ and, when equality holds, $j<K$.
Lemma~\ref{lem:decisive} gives $p\mid M_j$. For $j\ge1$,
\[
 \vp(M_j)+w_r(j)\ge1+\frac1{p-1}>w_r(K).
\]
For $j=0$, $M_0=0$ exactly. All remaining indices have weight at least $w_r(K)$ and
integral moments. At $K$ the moment is a unit. Thus the coefficient at $K$ realizes
the minimum, giving at least $K$ roots in the closed disk. The upper bound already
proved gives exactly $K$. Clearing denominators and applying the stability and
Eisenstein perturbation argument of Section~3 yields a polynomial with all supported
terms present, integer coefficients and distinct roots. This completes the proof
of Theorem~\ref{thm:sharp}.

\section{Consequences and checked examples}

\begin{corollary}[Exponential polynomials of arbitrary order]\label{cor:exp}
Among all nonzero sums \eqref{eq:exp} of total differential order at most $t$, with
frequencies in $\Z_p$, the greatest number of roots counted with multiplicity on
$\vp(z)>1/(p-1)$ is exactly $N_p^\circ(t)$. Likewise the closed-depth maximum equals
$N_p(t,r)$. In particular \eqref{eq:sharp} holds without requiring the amplitudes to
be constant or the frequencies to be ordinary integers.
\end{corollary}
\begin{proof}
Transport to an ordinary polynomial via Theorem~\ref{thm:transport} and use the
universal bound; the converse uses $P(\exp z)$ for a maximizing polynomial.
\end{proof}

\begin{corollary}[A degree-free obstruction to sparse annihilators]\label{cor:annihilator}
Let $\alpha$ be algebraic over a characteristic-zero number field $F$ embedded in
$\overline{\Q}_p$. If $u$ distinct $F$-conjugates of $\alpha$ lie in a common disk
$\D^\circ(a)$, every nonzero polynomial in $F[X]$ annihilating $\alpha$ and having
$t$ terms satisfies $u\le N_p^\circ(t)\le tM_{p,t}-1$ (for $t\ge2$).
In the range of Theorem~\ref{thm:sharp}, the sharper bound is
$u\le p^2+\lfloor(t-p)/2\rfloor-1$. A binomial or other lower-term obstruction is
obtained by padding to any allowed $t$. In particular $u>p^2+h-1$ rules out all
annihilators with at most $p+2h+1\le2p+1$ terms.
\end{corollary}
\begin{proof}
All $F$-conjugates are roots of an annihilator. Rescale the disk to center one and
apply the preceding theorems. A one-term annihilator has no nonzero root.
\end{proof}

This is directly applicable to the singular-modulus method of the original article:
any common deep-disk certificate with enough conjugates gives a degree-independent
sparsity exclusion. No new discriminant classification is asserted here, since that
would require the separate CM reduction and valuation computations. The coarse
inequality gives $t>\sqrt{u/p}$ when $t\ge2$; the residue-specific test may be much
stronger if the allowable exponent profiles are constrained.

\subsection{Exact finite checks}

The accompanying \texttt{verify.py} and \texttt{certificates.json} implement the
minimum-basis algorithm, construct the rational witnesses by Lagrange interpolation,
and independently count the roots of the resulting integer polynomials by evaluating
the Taylor coefficients $\sum_e c_e\binom ej$ and finding their last minimizing index.
Squarefreeness is checked by the exact polynomial gcd with the derivative. These
calculations are reproducible checks, not premises of the uniform proofs.

\begin{center}
\begin{tabular}{@{}rrrr@{}}
\toprule
$p$&$t$&$M_{p,t}$&$N_p^\circ(t)$\\
\midrule
2&2&1&1\\
2&3&1&2\\
2&4&2&4\\
2&5&2&4\\
2&6&2&8\\
2&7&2&8\\
2&8&4&9\\
2&9&4&16\\
\midrule
3&2&1&1\\
3&3&3&3\\
3&4&3&3\\
3&5&3&9\\
3&6&3&9\\
3&7&3&10\\
3&8&3&10\\
\bottomrule
\end{tabular}
\end{center}
All table entries were computed exhaustively over residue multisets using
Theorem~\ref{thm:algorithm}; the last several entries lie outside the explicit range
\eqref{eq:sharp}. For instance the $p=2,t=6$ maximum is $8$, showing that blindly
extrapolating \eqref{eq:sharp} past $2p+1$ would fail. The checks cover 15 complete
profile searches, 21 explicitly generated squarefree witnesses for primes through 19,
and 72 comparisons of different supports with identical residue data at several
depths (including the open disk). All passed; see \texttt{verification.log}.

\begin{center}
\begin{tabular}{@{}rrrr@{}}
\toprule
$p$ & $t=p+2$ & $t=2p+1$ & $N_p^\circ(2p+1)$\\
\midrule
2 &4&5&4\\
3 &5&7&10\\
5 &7&11&27\\
7 &9&15&52\\
11&13&23&126\\
101&103&203&10251\\
\bottomrule
\end{tabular}
\end{center}
The final table is the analytic formula, not an exhaustive search at $p=101$.

\section{Relation to the GRH proof and status of novelty}

\paragraph{Repository analysis.}
The root \texttt{README.md} of the specified directory reports 209 merged research
reports, the upstream proof of Dirichlet RH/GRH pinned at commit
\texttt{72fead102f663d5c7a407964a13dc0cf9d977207}, and the separate, unbuilt vendored
Lean development. The article's source section \cite{source} describes the transfer
architecture: arithmetic estimates for short integral forms, a spanning estimate,
and a determinant lower bound, followed by a one-sided analytic zero detector.
The article itself distinguishes unconditional results, Dirichlet-GRH consequences,
and reductions whose desired analytic bounds remain open. The singular-moduli section
uses integral moment recurrences, residue algebras, and Newton-polygon root counts;
the present result combines these with a determinant-minimizing integral lattice.
It is a transfer of proof ideas, not a logical consequence of GRH.

The repository snapshot used here was accessed through GitHub on 27 November 2026.
The principal inspected source blobs are listed below.
\begin{center}\small\begin{tabular}{@{}ll@{}}
\texttt{smd.tex}&\texttt{790993c6656d292056ac5eadc893061fa7a19002}\\
\texttt{src.tex}&\texttt{b0ac2147291b7da3a74cdf6072cc557b1e1f5ee8}\\
Article README&\texttt{b986fb54c6e641e9241439476bb5579e234fecd6}\\
\end{tabular}\end{center}
See the accompanying provenance note for other inspected sections.
The article is a merge of earlier reports, not an independent validation of all its
claims. Neither the full upstream Lean proof nor the 982-page merged article was
rechecked here; this paper is self-contained modulo the standard facts explicitly
used (monic polynomial division, Hermite interpolation, Legendre's formula,
$p$-adic exp/log and Newton polygon/Weierstrass preparation).

\paragraph{Prior literature.}
Van der Poorten and Rumely \cite{vpr} give bounds on zeros of $p$-adic exponential
polynomials in their domains of convergence; van der Poorten \cite{vdP} develops
Hermite-interpolation estimates. Rojas and Zhu \cite{rz} study algorithmic root
questions for very sparse polynomials over fixed local fields. These are related
but distinct from an exact transport of the full geometric root-depth multiset
between supports with identical finite residue data. A targeted literature search
was performed; it did not establish priority worldwide. The claimed advance is
precise relative to the inspected ProveIt article: arbitrary-order residue transport,
finite sharp determination at all depths, and a proof of the explicit new interval
\eqref{eq:sharp}. The open-ended task of obtaining compact closed forms for all
larger $t$ is not finished.

\paragraph{The prominent-conjecture target.}
Another investigated route was the article's all-degrees symmetric-group
positive-monomial-character question (\texttt{nab:q:un}, \texttt{grp:q:vn}), which
would sharpen a zero-multiplicity bound for Dedekind zeta functions. No uniform
positive decomposition was established here. We do not claim to solve that question,
Artin's holomorphy conjecture, nonabelian GRH, the six-term singular-modulus
classification, or another famous conjecture. The unconditional exact theorems above
are the completed mathematical contribution. They have not yet been formalized in
Lean or independently refereed.

\begin{thebibliography}{9}
\bibitem{article} ProveIt, \emph{Beyond the Riemann hypothesis for Dirichlet
L-functions: consequences, transfers and limits}, singular-moduli section,
\href{https://github.com/VladimirReshetnikov/ProveIt/blob/main/Analysis/GeneralizedRiemannHypothesis/Article/sections/smd.tex}{\texttt{Article/sections/smd.tex}},
Theorems \texttt{smd:thm:rankrule}, \texttt{smd:thm:jump}, and Question \texttt{smd:q:roots}.
Blob SHA: \texttt{790993c6656d292056ac5eadc893061fa7a19002}.
\bibitem{source} ProveIt, source theorem and transferable mechanism,
\href{https://github.com/VladimirReshetnikov/ProveIt/blob/main/Analysis/GeneralizedRiemannHypothesis/Article/sections/src.tex}{\texttt{Article/sections/src.tex}}.
\bibitem{vpr} A. J. van der Poorten and R. S. Rumely, Zeros of $p$-adic exponential
polynomials II, \emph{J. London Math. Soc.} (2) \textbf{36} (1987), 1--15.
\href{https://doi.org/10.1112/jlms/s2-36.1.1}{doi:10.1112/jlms/s2-36.1.1}.
\bibitem{vdP} A. J. van der Poorten, Hermite interpolation and $p$-adic exponential
polynomials, \emph{J. Austral. Math. Soc. Ser. A} \textbf{22} (1976), 11--23.
\href{https://www.cambridge.org/core/journals/journal-of-the-australian-mathematical-society/article/hermite-interpolation-and-padic-exponential-polynomials/46C55A4B6B0E76934DBD8CECFA35749A}{Publisher page}.
\bibitem{rz} J. M. Rojas and Y. Zhu, Root repulsion and faster solving for very sparse
polynomials over $p$-adic fields, \emph{J. Number Theory} \textbf{241} (2022), 655--699.
\href{https://www.sciencedirect.com/science/article/abs/pii/S0022314X22000877}{Publisher page}.
\end{thebibliography}
\end{document}
